\scanpage{255}
\mspage{129}
% archive sheet number: 254

Pour prouver que
\[
 f_0=f_1\quad\text{dans }\mathrm{Hom}_W(X,Y),
\]
il suffit de voir que \deleted{$\partial'_0{}^X$ et $\partial^{\unsure{''|'}}_1$}
\[
 \partial'_0{}^X=\partial'_1{}^X\quad\text{dans }\mathrm{Hom}_W(X,J\times X).
\]

Or en fait $\partial'_0{}^X$ et $\partial'_1{}^X$ sont homotopes, car
\diagram{\begin{tikzcd}[ampersand replacement=\&,column sep=small]
 e\arrow[r,shift left=.6ex,"j_0"]\arrow[r,shift right=.6ex,"j_1"'] \& J
\end{tikzcd}}
sont homotopes, \deleted{\unsure{d'un}} ($j_0$ et $j_1$ \unsure{appartiennent|appart.} à la $=$ comp. conn. de $J$), donc aussi \unsure{leur produit|leurs produits} par $\mathrm{id}_X$.

On voit donc que pour deux objets quelconques $X,Y$ de \deleted{$\mathrm{Cat}$, $\unsure{\mathrm{Hot}}$} $\mathrm{Cat}$, \unsure{toutes} les flèches de $X$ dans $Y$ sont égales dans $\mathrm{Hot}$. L'argument fait prouve $=$ que \deleted{toutes} les flèches \uline{dans $\mathrm{Hot}_W$} de $X$ dans $Y$ sont égales.

Cela signifie que \unsure{la} cat. $\mathrm{Hot}^{\unsure{*|\bullet}}_W$ (où \inserted{pour} deux objets quelc. \inserted{$X,Y$}, \deleted{\unsure{sont reliés par une flèche} \inserted{$\mathrm{Hom}(X,Y)$}} $\mathrm{Hom}_{\mathrm{Cat}}(X,Y)$ et \unsure{a|à} fortiori $\mathrm{Hom}_W(X,Y)$ est $\ne\varnothing$) est équivalente à la cat. ponctuelle.

Quant à $\mathrm{Hot}_W$ elle $=$, \deleted{qui est}\inserted{dont l'ens} des objets est $\{\varnothing\}\amalg\mathrm{Ob}\,\mathrm{Hot}^{*}_W$, \inserted{avec $\varnothing$ objet initial}, elle sera donc équivalente à la catégorie ponctuelle, ou bien à $\Delta^1$, selon que $\varnothing\to e$ est un iso dans $\mathrm{Hot}_W$, ou non.
